\documentclass[a4paper,12pt]{article}

\title{Blocco 3}
\author{Leonardo Centazzo}
\date{\today}


% Pacchetti essenziali
\usepackage{amsmath, amssymb, amsthm}
\usepackage{enumitem}
\usepackage{geometry}

% Impostazioni pagina minime
\geometry{a4paper, margin=2cm}

% Stili per teoremi, proposizioni, ecc.
\theoremstyle{plain}
\newtheorem{teorema}{Teorema}
\newtheorem{proposizione}{Proposizione}
\newtheorem{lemma}{Lemma}
\newtheorem{corollario}{Corollario}

\theoremstyle{definition}
\newtheorem{definizione}[teorema]{Definizione}
\newtheorem{esempio}[teorema]{Esempio}
\newtheorem{esercizio}{Esercizio}
\newtheorem{osservazione}[teorema]{Osservazione}

\newtheorem*{warning}{Warning}

\theoremstyle{remark}
\newtheorem{nota}[teorema]{Nota}


\begin{document}

\maketitle

For the next two problems, suppose that $\mathcal{F}$ is a nonprincipal, $k$-complete filter on $k$.
Define the dual ideal $I = \{A \subseteq k : k \setminus A \in \mathcal{F}\}$.
We say that $A \subset k$ is a positive set if $A \not \in I$ (equivalently, for all $C \in \mathcal{F}$, $C \cap A \not = \emptyset$).
For a positive set $A$ we say that $f: A \to k$ is regressive if for all $\alpha \in A$, $f(\alpha) < \alpha$.

\begin{esercizio}

	Suppose that for all $A \subseteq k$, $A \not \in I$ and regressive function $f: A \to k$, $f$ is constant on a positive subset of $A$.
	Show that $\mathcal{F}$ is closed under diagonal intersections.

\end{esercizio}

\begin{proof}

	Let $C_i \in \mathcal{F}$ for $i < k$ and let $D = \triangle_i C_i$. 
	By contradiction, assume $D \not \in \mathcal{F}$. Then $D^{c} \not \in I$, which means that $D^{c}$ is positive.
	Consider $f: D^{c} \to k, \alpha \mapsto \text{min}\{i < \alpha : \alpha \not \in C_i \}$.
	\begin{itemize}
		\item $f$ well defined:
			$\alpha \in D^{c} \Rightarrow \exists i < \alpha \ \big( \alpha \not \in C_i \big)$, so we have that 
			for all $\alpha \in D^{c}$, $\{i < \alpha : \alpha \not \in C_i\} \not = \emptyset$.
			It has minimum because $\subseteq ON$.

		\item $f$ regressive: by definition
	\end{itemize}

	Let $P \subseteq D^{c}$ positive such that $f \restriction{P} = cost_{\beta}$.
	Now observe that $P \cap C_{\beta} = \emptyset$: in fact $p \in P \Rightarrow f(p) = \beta \Rightarrow p \not \in C_{\beta}$.
	Thus we reached a contradiction because $P$ is positive, $C_{\beta} \in \mathcal{F}$ but $P \cap \mathcal{F} = \emptyset$.

\end{proof}





\begin{esercizio}

	Suppose that $\mathcal{F}$ is closed under diagonal intersections.
	Show that for all $A \subseteq k, A \not \in I$ and regressive $f: A \to k$, $f$ is constant on a positive subset of $A$.

\end{esercizio}

\begin{proof}
	For all $i < k$, let $N_i = f^{-1}\{i\}$. By contradiction, assume for all $i < k$, $N_i$ is negative (meaning not positive).
	Then for all $i < k$, let $C_i \in \mathcal{F}$ such that $C_i \cap N_i = \emptyset$.
	Consider $C = \triangle_i C_i$. Then $C \in \mathcal{F}$. Since $A$ is positive, we have $A \cap C \not = \emptyset$.
	Let $\alpha \in A \cap C$. Since $\alpha \in C$ we have that $\forall i< \alpha \big(\alpha \in C_i \big)$.
	This means that $\forall i<\alpha \big( f(\alpha) \not = i \big)$. Equivalently, $f(\alpha) \geq \alpha$, which is a contradiction
	since $f$ is regressive.

\end{proof}



For the next two problems let $k$ be a measurable cardinal with normal measure $U$ and let $P$ be the Prickry forcing at $k$.

\begin{esercizio}

	Suppose that $G$ is $P$-generic over the ground model $V$.
	Let $\vec{\alpha} = \bigcup_{(s,A) \in G} s$ and let 
	$G_{\vec{\alpha}} = \{(s,A) : \exists n \big( \vec{\alpha} \restriction n = s , \forall k \geq n \alpha_k \in A \big) \}$.
	Prove that $G = G_{\vec{\alpha}}$.

\end{esercizio}

\begin{proof}

	
	\begin{itemize}
		\item $G \subseteq G_{\vec{\alpha}}$:
			Let $(s,A) \in G$. Then by definition of $\vec{\alpha}$ we have that $s$ must be an initial segment of $\vec{\alpha}$.
			By contradiction, assume $\{n> |s| : \vec{\alpha}(n) \not \in A\} \not = \emptyset$.
			Let $m$ be minimum of the above set. Let $(t,B) \in G$ such that $|t| \geq m$.
			Such an element exists by genericity of $G$.
			Since $(s,A), (t,B)$ must be compatible we have that $t(n) \in A$ for all $n > |s|$. 
			In particular $\vec{\alpha}(m) = t(m) \in A$, contradiction.
			
		\item $G_{\vec{\alpha}} \subseteq G$:
			Let $(s,A) \in G_{\vec{\alpha}}$. Consider $D = \{(t,B) : |t| \geq |s|, B \subseteq A\}$.
			$D$ is dense because if we take $(z,C) \in P$ (if necessary) we can extend $z$ with the least elements of $C$ until
			we reach length $|s|$. Let's say we concatenated $\overline{c}$. 
			Then $(z,C) \geq (z \frown \overline{c}, C \setminus \overline{c}) 
			\geq (z \frown \overline{c}, (C \setminus \overline{c}) \cap A) \in D$.
			Since $G$ generic and $D$ dense, we have $G \cap D \not = \emptyset$. 
			Let $(t,B) \in G \cap D$. From $(t,B) \in G$ we get $t$ initial segment of $\vec{\alpha}$.
			Combined with $(t,B) \in D$ we deduce $t$ extends $s$ with elements in $A$.
			Finally from $(t,B) \in D$ we deduce $B \subseteq A$ and conclude $(t,B) \leq (s,A)$.
			Since $(t,B) \in G$ we have $(s,A) \in G$.

	\end{itemize}
\end{proof}



All the notation for the following exercise is intended as the one used in the lessons unless otherwise specified.
So for example $j: V \to M$ is elementary embedding, $M$ is the ultraproduct, $k^{M} = [Id]$, $M^{k} \subseteq M$, $M \subseteq V$.

\begin{esercizio}
	
	Suppose that in the ground model $V$, $2^{k} = k^{++}$ and $G$ is $P$-generic over $V$.
	Suppose that $\langle k_n : n < \omega \rangle$ is the Prikry sequence added by $G$.
	Show that for all large $n$, $k_n$ is a cardinal and $2^{k_n} = k_n ^{++}$.
	
\end{esercizio}

\begin{proof}		

	Consider $B = \{\lambda \in k : \lambda \text{ is a cardinal and } 2^{\lambda} = \lambda^{++}\}$.
	We claim that $B \in U$. This is equivalent to $M \models 2^{k} = k^{++}$ by Los Theorem.
	\begin{itemize}
		\item $(k^{+})^M = (k^{+})^V$ : In fact, since $M \subseteq V$ we have $(k^{+})^M \leq (k^{+})^{V}$. 
			If by contradiction $(k^{+})^{M} < (k^{+})^{V}$, then there would exist $f \in V, f:k \to (k^{+})^{M}$ surjective.
			Since $M^{k} \subseteq M$ we have $f \in M$, so $(k = k^{+})^M$, a contradiction.
			In particular from now on we simply write $\lambda$ to denote $(k^{+})^{V}, (k^{+})^{M}$, since they are the same.

		\item $(2^{k})^M = (2^{k})^V$: Since $M \subseteq V$ we have $(2^{k})^M \subseteq (2^{k})^V$. Since $M^{k} \subseteq M$ the other
			inclusion holds as well.

		\item $(k^{++})^M = (k^{++})^V$: I don't know if it is true. In any case I'd like to receive the solution of the exercise.

	\end{itemize}
	Now since in $V$ it is true that $2^{k} = k^{++}$ and all cardinals in the equation are the same in $M$, we deduce that the equation
	holds in $M$ as well.

	From $B \in U$ we deduce that $D = \{(s,A) : A \subseteq B\}$ is dense.
	Then $G \cap D \not = \emptyset$. Let $(s,A) \in G \cap D$.
	Then for all $n > |s| \ \big( k_n \in A \big)$, which implies for all $n > |s| \ \big( k_n \text{ is a cardinal and } 2^{k_n} = k_n ^{++}\big)$.

\end{proof}


\end{document}

